\documentclass[12pt]{article}
\usepackage[utf8]{inputenc}
\usepackage[T1]{fontenc}
\usepackage[margin=1in]{geometry}
\usepackage{setspace}
\usepackage{xcolor}
\usepackage[colorlinks=true, urlcolor=blue!60!black, linkcolor=blue!60!black]{hyperref}
\setlength{\parindent}{2em}
\setlength{\parskip}{1.25em}
\renewcommand{\baselinestretch}{1.0}
\begin{document}
\begin{center}
{\LARGE \textit{Nixon v. Administrator of General Services}, 433 U.S. 425 (1977)\par}
\vspace{0.8em}
{\large Case Brief --- Four-Page Edition\par}
\end{center}
\vspace{1em}

\section{Memory Jogger}

After Watergate, Congress directed the Administrator of General Services to take custody of former President Nixon's papers and tapes for screening by executive-branch archivists; the Court held 7--2 that the Act was not a bill of attainder (Nixon was a legitimate ``class of one'' and was not punished) and did not violate separation of powers, presidential privilege, privacy, or association rights (\href{https://supreme.justia.com/cases/federal/us/433/425/}{Justia}).

\section{Detailed Case Facts}

President Nixon resigned on August 9, 1974, after the House Judiciary Committee approved articles of impeachment. On September 8, the day President Ford pardoned him, GSA Administrator Arthur Sampson announced an agreement reciting that Nixon retained title to some 42 million pages and 880 tapes. Nixon could withdraw documents after three years and order tapes destroyed after September 1, 1979, and all tapes would be destroyed at his death or by 1984.

Congress responded with the Presidential Recordings and Materials Preservation Act, signed December 19, 1974. Title I directed the Administrator to take complete possession of Nixon's presidential materials, barred their destruction, made them available for judicial process subject to privileges, guaranteed Nixon access, and required regulations for archival screening that would protect privileges and return purely personal materials. Section 105(c) provided just compensation for any taking, and Title II created a commission to study records of future Presidents.

\section{Procedural History}

Nixon sued the next day on facial constitutional grounds. Because public-access regulations had not taken effect, a three-judge district court reviewed only the custody and screening provisions, rejected every claim, and dismissed. The Supreme Court noted probable jurisdiction, heard argument on April 20, 1977, and affirmed on June 28, 1977.

\section{Judicial Votes}

Affirmed, \textbf{7--2}.
\begin{itemize}
\item \textbf{Majority:} Justice Brennan, joined by Justices Stewart, Marshall, and Stevens; Justice White in all but Part VII; Justice Powell in all but Parts IV and V; and Justice Blackmun as to Part VII.
\item \textbf{Concurring:} Justice Stevens; Justices White, Blackmun, and Powell concurred in part and in the judgment.
\item \textbf{Dissenting:} Chief Justice Burger and Justice Rehnquist.
\end{itemize}

\section{Holding}

Title I is facially valid. It is not a bill of attainder: Nixon was a ``legitimate class of one,'' and the Act inflicted no punishment under historical, functional, and motivational tests. It does not violate separation of powers, because it does not unduly prevent the Executive from performing its functions, particularly since custody stayed within the Executive Branch and Presidents Ford and Carter supported it. A former President may assert the qualified presidential privilege, but his claim carries diminished weight without the incumbent's support and yields to minimal archival intrusion. Nixon's privacy and associational interests were outweighed by the public interest and protected by screening and return provisions.

\section{Analysis of Opinions}

\textbf{Majority (Brennan).} \emph{Attainder:} \textit{United States v. Brown} does not bar every law singling out an individual. Every President from Hoover through Johnson had deposited papers in libraries; only Nixon's agreement threatened destruction. There was no traditional punishment, compensation negated confiscation, the Act served nonpunitive goals of preserving evidence and history, and the legislative record showed regulatory concern, not condemnation. \emph{Separation of powers:} rejecting an ``archaic view of the separation of powers as requiring three air-tight departments of government,'' the Court asked whether the Act prevents the Executive from accomplishing its functions and, if so, whether an overriding need justifies it. GSA archivists are executive employees, and Congress received no privileged access. \emph{Privilege:} confidentiality survives a President's tenure, but screening by discreet professionals resembles the \textit{in camera} review approved in \textit{United States v. Nixon}. \emph{Privacy:} under \textit{Whalen v. Roe}, personal materials were inextricably commingled with official ones, so screening was the least intrusive way to separate them.

\textbf{Concurrences.} Stevens stressed that Nixon's resignation and pardon made him unique and warned against laws targeting an individual occupant. White insisted that private diaries be returned promptly. Blackmun warned against using the Act as a template in hostile transitions. Powell treated Ford's and Carter's support as dispositive and warned that ``difficult constitutional questions lie ahead.''

\textbf{Dissents.} Burger saw principles bent under ``hydraulic pressure'' and an unconstitutional attainder. Rehnquist warned that the ruling would ``daily stand as a veritable sword of Damocles over every succeeding President and his advisers.''

\section{Examples: Future Applications}

\begin{enumerate}
\item \textbf{Same-Side (No Attainder): One-Company Procurement Ban.} Barring agencies from contracting with a named vendor with documented security vulnerabilities is prophylactic regulation; the vendor may be a legitimate class of one.
\item \textbf{Same-Side (No Attainder): Securing an Official's Server.} A statute directing archivists to take custody of a former secretary's private email server, screen it, return personal messages, and pay compensation tracks the PRMPA.
\item \textbf{Opposite-Side (Attainder): Uncompensated Forfeiture.} Forfeiting the vendor's patents without compensation, after floor speeches calling its founders ``traitors,'' fails the historical and motivational tests, as in \textit{Foretich}.
\item \textbf{Opposite-Side (Separation of Powers): Records to a House Committee.} Giving a committee custody and publication power over privileged records, over the incumbent's objection, removes every feature that saved the PRMPA.
\item \textbf{Fence-Sitter: A Hostile Successor's Waiver.} An incumbent's waiver favors disclosure under \textit{Nixon}, but publication without any evidentiary need, and Blackmun's hostile-transition warning, make the outcome uncertain.
\end{enumerate}

\section{Critique}

\textbf{Retreat from \textit{Brown}.} Critics note that the majority drew on Justice White's \textit{Brown} dissent to turn a structural bar on targeted legislation into a deferential inquiry in which almost any plausible nonpunitive rationale rescues a named-target law.

\textbf{Pretextual classes and functionalist drift.} As Burger and Stevens warned, pairing a targeted statute with a forward-looking study commission risks normalizing partisan targeting, and relying on the incumbent's consent undervalues incentives in polarized transitions, as Blackmun foresaw.

\section{Key Quotations}

\begin{quote}
``The Court, therefore, often has looked beyond mere historical experience and has applied a functional test of the existence of punishment, analyzing whether the law under challenge, viewed in terms of the type and severity of burdens imposed, reasonably can be said to further nonpunitive legislative purposes.''
\end{quote}
\noindent --- Justice Brennan, for the Court, 433 U.S. at 475--476.

\begin{quote}
``Rather, in determining whether the Act disrupts the proper balance between the coordinate branches, the proper inquiry focuses on the extent to which it prevents the Executive Branch from accomplishing its constitutionally assigned functions. \textit{United States v. Nixon}, 418 U.S. at 711--712. Only where the potential for disruption is present must we then determine whether that impact is justified by an overriding need to promote objectives within the constitutional authority of Congress.''
\end{quote}
\noindent --- Justice Brennan, for the Court, 433 U.S. at 443.

\begin{quote}
``In short, while the Bill of Attainder Clause serves as an important `bulwark against tyranny,' \textit{United States v. Brown}, 381 U.S. at 443, it does not do so by limiting Congress to the choice of legislating for the universe, or legislating only benefits, or not legislating at all.''
\end{quote}
\noindent --- Justice Brennan, for the Court, 433 U.S. at 471.

\end{document}
